\documentclass[12pt,letterpaper,oneside]{report}

% Paquetes de idioma y codificación
\usepackage[utf8]{inputenc}
\usepackage[T1]{fontenc}
\usepackage[spanish,es-tabla]{babel}

% Paquetes de diseño y formato
\usepackage[margin=2.5cm, headheight=15pt]{geometry}
\usepackage{listings}
\usepackage{xcolor}
\usepackage{hyperref}
\usepackage{titlesec}
\usepackage{graphicx}
\usepackage{amsmath}
\usepackage{amssymb}
\usepackage{booktabs}
\usepackage{fancyhdr}
\usepackage{microtype}
\usepackage{tcolorbox}
\usepackage{makeidx}
\usepackage{caption}
\usepackage{enumitem}
\usepackage{longtable}
\makeindex

% Colores personalizados para el resaltado
\definecolor{yamlkey}{RGB}{0, 102, 204}
\definecolor{yamlval}{RGB}{204, 0, 0}
\definecolor{yamlcomment}{RGB}{120, 120, 120}
\definecolor{bglight}{RGB}{245, 245, 245}
\definecolor{darkgray}{RGB}{50, 50, 50}
\definecolor{pythonkw}{RGB}{0, 0, 150}
\definecolor{pythonstr}{RGB}{0, 150, 0}
\definecolor{notebackground}{RGB}{240, 248, 255}
\definecolor{noteborder}{RGB}{0, 102, 204}
\definecolor{warnbackground}{RGB}{255, 245, 230}
\definecolor{warnborder}{RGB}{204, 102, 0}

% Configuración de enlaces
\hypersetup{
    colorlinks=true,
    linkcolor=yamlkey,
    filecolor=yamlval,
    urlcolor=yamlkey,
    citecolor=yamlkey,
    pdftitle={Manual de Usuario - Fenix FEM},
    pdfauthor={Jaime Retama Velasco}
}

% Formato de títulos
\titleformat{\chapter}[display]
  {\normalfont\huge\bfseries\color{yamlkey}}{\chaptertitlename\ \thechapter}{20pt}{\Huge}
\titleformat{\section}{\Large\bfseries\color{darkgray}}{\thesection}{1em}{}
\titleformat{\subsection}{\large\bfseries\color{darkgray}}{\thesubsection}{1em}{}
\titleformat{\subsubsection}{\large\bfseries\color{darkgray}}{\thesubsubsection}{1em}{}

% Encabezados y pies de página
\pagestyle{fancy}
\fancyhf{}
\fancyhead[L]{\textbf{\color{darkgray}Fenix FEM}}
\fancyhead[R]{\color{darkgray}\leftmark}
\fancyfoot[C]{\thepage}
\renewcommand{\headrulewidth}{0.4pt}
\renewcommand{\footrulewidth}{0.4pt}

% Cajas de notas y advertencias
\newtcolorbox{nota}[1][]{colback=notebackground,colframe=noteborder,title=\textbf{Nota},#1}
\newtcolorbox{advertencia}[1][]{colback=warnbackground,colframe=warnborder,title=\textbf{Advertencia},#1}

% Lenguaje YAML para listings
\lstdefinelanguage{yaml}{
  keywords={nodes, materials, elements, mesh, mesh_material, mesh_thickness, mesh_quadrature, mesh_physical_groups, quadrature, boundary_conditions, boundary_conditions_by_node, boundary_conditions_by_coord, boundary_conditions_by_group, point_loads, point_loads_by_node, point_loads_by_coord, point_loads_by_group, solver, output, file_name, pre_process, export, results, text_export, frequency, nodal_results, element_results, format, id, type, coords, ref_vector, A, I, Iy, Iz, J, As, E, nu, sigma_y, H, hypothesis, ux, uy, uz, rx, ry, rz, tol, coord, val, num_steps, max_iter, adaptive, tol_iter, max_lambda, initial_dl, max_steps, material, thickness, kappa_0, alpha},
  keywordstyle=\color{yamlkey}\bfseries,
  ndkeywords={true, false, null, all, plane_stress, plane_strain, LinearSolver, NonlinearSolver, ArcLengthSolver, Elastic1D, Elastic2D, Elastoplastic1D, VonMises2D, IsotropicDamage1D, IsotropicDamage2D, CableMaterial1D, Truss2D, Truss2DCorot, Truss3D, Truss3DCorot, Cable2DCorot, Cable3DCorot, Frame2DEuler, Frame2DTimoshenko, Frame2DEulerCorot, Frame3D, Quad4, Tri3},
  ndkeywordstyle=\color{yamlval}\bfseries,
  identifierstyle=\color{black},
  sensitive=false,
  comment=[l]{\#},
  commentstyle=\color{yamlcomment}\ttfamily,
  stringstyle=\color{yamlval}\ttfamily,
  morestring=[b]',
  morestring=[b]"
}

% Configuración general de listings
\lstset{
  basicstyle=\ttfamily\footnotesize,
  backgroundcolor=\color{bglight},
  frame=single,
  rulecolor=\color{lightgray},
  breaklines=true,
  showstringspaces=false,
  extendedchars=true,
  captionpos=b,
  literate={á}{{\'a}}1 {é}{{\'e}}1 {í}{{\'i}}1 {ó}{{\'o}}1 {ú}{{\'u}}1
           {Á}{{\'A}}1 {É}{{\'E}}1 {Í}{{\'I}}1 {Ó}{{\'O}}1 {Ú}{{\'U}}1
           {ñ}{{\~n}}1 {Ñ}{{\~N}}1 {ü}{{\"u}}1 {Ü}{{\"U}}1 {¿}{{?`}}1 {¡}{{!`}}1
}

\begin{document}

% --- PORTADA ---
\begin{titlepage}
    \centering
    \vspace*{2cm}
    {\Huge\bfseries\color{yamlkey} Fenix FEM \par}
    \vspace{1cm}
    {\LARGE Plataforma de Análisis por Elementos Finitos en Python \par}
    \vspace{0.5cm}
    \rule{\linewidth}{0.5mm} \par
    \vspace{2cm}
    {\huge\bfseries Manual de Usuario y Guía de Referencia \par}
    \vspace{2cm}
    {\Large Documentación Oficial y Arquitectura del Sistema \par}
    \vspace{0.5cm}
    {\large Versión 2.0 \par}
    \vspace{1.5cm}
    {\Large \textbf{Autor:} Jaime Retama Velasco \par}
    \vspace{0.5cm}
    {\large Facultad de Estudios Superiores Aragón \\ Universidad Nacional Autónoma de México \par}
    \vfill
    {\large \today \par}
\end{titlepage}

\pagenumbering{roman}
\setcounter{page}{1}

\tableofcontents
\newpage

\chapter*{Resumen Ejecutivo}
\addcontentsline{toc}{chapter}{Resumen Ejecutivo}
\noindent \textbf{Fenix FEM} es un programa de elementos finitos en aproximación de desplazamientos, escrito en Python, dirigido a la investigación en mecánica de sólidos y al análisis estructural. Cubre problemas mecánicos en régimen lineal y no lineal --- tanto material (plasticidad, daño, elasticidad unilateral) como geométrico (formulaciones corotacionales) --- y está pensado para crecer en el tiempo incorporando nuevos elementos, materiales y métodos de solución sin reescribir el núcleo del programa.

El catálogo actual incluye:

\begin{itemize}
    \item \textbf{10 elementos finitos 1D} en el plano y en el espacio: armaduras, cables y marcos en variantes lineales y corotacionales.
    \item \textbf{2 elementos finitos 2D} continuos: cuadrilátero bilineal \texttt{Quad4} y triángulo lineal \texttt{Tri3}.
    \item \textbf{7 modelos constitutivos}: elasticidad lineal 1D y 2D, plasticidad J2 1D y 2D, daño isótropo 1D y 2D, y elasticidad unilateral para cables.
    \item \textbf{3 esquemas de solución}: solver lineal directo, Newton-Raphson incremental con paso adaptativo, y longitud de arco cilíndrico de Crisfield para snap-through y snap-back.
\end{itemize}

La interacción del usuario con el motor se realiza mediante un único \textbf{archivo de entrada YAML} que describe geometría, materiales, condiciones de frontera, cargas y configuración del solver. La salida se exporta en formato \texttt{.vtu} (VTK XML) para post-procesamiento en \textit{ParaView}, y opcionalmente en texto plano estilo FEAP.

Este manual documenta el uso del programa --- sintaxis del archivo YAML, catálogos de componentes y workflow de ejecución. Para detalles internos sobre arquitectura, formulaciones físicas y workflow de extensión por desarrollo asistido por IA, consulte la documentación interna del repositorio (\texttt{docs/}, \texttt{Reglas.md}).

\newpage

\pagenumbering{arabic}
\setcounter{page}{1}

% ==========================================
\chapter{Introducción y Configuración del Entorno}
% ==========================================

\section{Filosofía del Proyecto}
Fenix FEM nace como herramienta de investigación en mecánica del medio continuo y análisis no lineal de estructuras. Su diseño persigue dos objetivos simultáneos:

\begin{enumerate}
    \item \textbf{Rigor numérico y físico.} Cada formulación incorporada está validada contra solución analítica o benchmark establecido. La separación estricta \texttt{Elemento} $\leftrightarrow$ \texttt{Material} permite combinar cualquier cinemática con cualquier ley constitutiva siempre que sus protocolos sean compatibles.
    \item \textbf{Crecimiento sin fricción.} La arquitectura está optimizada para que añadir un elemento, material o solver nuevo sea barato. Auto-registro vía decoradores, contratos declarativos (\texttt{STRAIN\_DIM}, \texttt{DOF\_NAMES}) y validación temprana del input son los mecanismos centrales.
\end{enumerate}

\subsection{Modelo \textit{Data-Driven}}
Toda la orquestación de un análisis se concentra en un archivo \texttt{.yaml} legible por humanos. El usuario no necesita escribir código Python para definir un modelo; la API programática queda reservada para casos avanzados (optimización paramétrica, mallas generadas proceduralmente, acoplamiento con otros sistemas).

\section{Requisitos del Sistema}\index{Instalación}
Fenix FEM se ejecuta sobre Python 3.8 o superior en Windows, Linux y macOS sin compilación nativa.

\subsection{Dependencias}
\begin{lstlisting}[language=bash]
pip install numpy scipy pyyaml meshio numba
\end{lstlisting}

\begin{itemize}
    \item \textbf{\texttt{numpy} / \texttt{scipy}}: álgebra lineal densa y dispersa; \texttt{spsolve} sobre matrices CSR.
    \item \textbf{\texttt{pyyaml}}: análisis sintáctico del archivo de entrada.
    \item \textbf{\texttt{meshio}}: importación de mallas Gmsh y exportación a VTK.
    \item \textbf{\texttt{numba}}: compilación JIT de los kernels críticos en elementos 2D y plasticidad.
\end{itemize}

\section{Herramientas Externas Recomendadas}
\begin{itemize}
    \item \textbf{Gmsh}\index{Gmsh}: generador de mallas; produce archivos \texttt{.msh} consumibles directamente desde el YAML.
    \item \textbf{ParaView}\index{ParaView}: visualización interactiva de los archivos \texttt{.vtu} generados.
\end{itemize}

\section{Ejecución de un Análisis}
El runner principal vive en \texttt{examples/ejecutar\_yaml.py}. Para ejecutar un modelo:

\begin{lstlisting}[language=bash]
python examples/ejecutar_yaml.py ruta/al/modelo.yaml
\end{lstlisting}

El script carga el YAML, valida su contenido (acumulando \emph{todos} los errores en una sola pasada antes de abortar), construye el dominio, ensambla el solver pedido, ejecuta el análisis y exporta los resultados según el bloque \texttt{output}.

% ==========================================
\chapter{Arquitectura del Programa}
% ==========================================

\section{Separación Elemento $\leftrightarrow$ Material}
Cada elemento finito declara dos contratos sobre el material que acepta:

\begin{itemize}
    \item \textbf{\texttt{STRAIN\_DIM}}: dimensión Voigt de la deformación que entrega al material. \texttt{1} para elementos axiales (truss, cable, marcos --- la fibra centroidal); \texttt{3} para 2D (\texttt{[$\varepsilon_{xx}, \varepsilon_{yy}, \gamma_{xy}$]}); \texttt{6} para 3D continuo (no implementado todavía).
    \item \textbf{\texttt{DOF\_NAMES}}: lista de grados de libertad por nodo. Por ejemplo, \texttt{['ux', 'uy', 'rz']} para un marco 2D.
\end{itemize}

El protocolo del material es uniforme:
\begin{lstlisting}[language=Python]
sigma, E_tangent, new_state = material.compute_state(epsilon, state_vars)
\end{lstlisting}

donde \texttt{epsilon} tiene dimensión \texttt{STRAIN\_DIM}, \texttt{sigma} es el esfuerzo en la misma representación, \texttt{E\_tangent} es la matriz tangente consistente y \texttt{state\_vars} encapsula las variables internas (deformación plástica, daño acumulado, etc.).

Un material declara igualmente su \texttt{STRAIN\_DIM}; el dominio rechaza al construir cualquier emparejamiento incompatible.

\section{Auto-Registro de Componentes}
Elementos, materiales y solvers se registran automáticamente al arrancar el programa mediante decoradores:

\begin{lstlisting}[language=Python]
@ElementRegistry.register
class Truss2D(Element):
    DOF_NAMES = ['ux', 'uy']
    STRAIN_DIM = 1
    ...
\end{lstlisting}

Esto permite que el archivo YAML los referencie por nombre (\texttt{type: Truss2D}) sin que el usuario tenga que importar nada manualmente. El módulo \texttt{fenix.autodiscover} recorre los paquetes \texttt{elements/}, \texttt{materials/} y \texttt{math/solvers.py} en la primera importación y puebla los registries.

\section{Validación Temprana del Input}
El parser YAML acumula todos los errores detectados en una sola pasada (no aborta al primero) y los presenta juntos. Comprueba:

\begin{itemize}
    \item Ausencia de bloques obligatorios (\texttt{nodes}, \texttt{materials}, \texttt{elements}, \texttt{solver}).
    \item IDs duplicados de nodo, material o elemento.
    \item Referencias a materiales o nodos inexistentes desde \texttt{elements} y \texttt{boundary\_conditions}.
    \item Tipos de elemento, material y solver no registrados.
    \item \textbf{Parámetros no aceptados por el constructor del elemento.} Si un elemento define \texttt{A} e \texttt{I} como parámetros, escribir \texttt{large\_strains: true} en su entrada YAML produce un error explícito que indica los parámetros válidos para ese tipo.
\end{itemize}

% ==========================================
\chapter{Archivo de Entrada YAML}\index{YAML}
% ==========================================

La interacción exclusiva con el motor de Fenix FEM se realiza mediante el archivo de entrada \texttt{.yaml}. Un archivo válido contiene los bloques siguientes (algunos opcionales según el caso de uso):

\begin{itemize}
    \item \texttt{nodes} \emph{o} \texttt{mesh} --- geometría.
    \item \texttt{materials} --- catálogo de leyes constitutivas instanciadas.
    \item \texttt{elements} --- conectividad (omisible si se importa malla Gmsh).
    \item \texttt{boundary\_conditions\_by\_node} / \texttt{\_by\_coord} / \texttt{\_by\_group} --- restricciones de Dirichlet.
    \item \texttt{point\_loads\_by\_node} / \texttt{\_by\_coord} / \texttt{\_by\_group} --- cargas de Neumann puntuales.
    \item \texttt{solver} --- tipo de solver y sus parámetros.
    \item \texttt{output} --- frecuencia y formato de exportación de resultados.
\end{itemize}

\section{Geometría: \texttt{nodes} y \texttt{elements}}

Para modelos pequeños o estructuras reticulares (armaduras, marcos, cables), la geometría se escribe directamente en el YAML:

\begin{lstlisting}[language=yaml]
nodes:
  - {id: 1, coords: [0.0, 0.0]}
  - {id: 2, coords: [2.0, 0.0]}
  - {id: 3, coords: [4.0, 0.0]}

elements:
  - id: 1
    type: Frame2DEuler
    material: 1
    nodes: [1, 2]
    A: 0.01
    I: 8.33e-6
\end{lstlisting}

Las coordenadas pueden ser de 2 o 3 componentes; los elementos 1D 3D (\texttt{Truss3D}, \texttt{Cable3DCorot}, \texttt{Frame3D}) requieren las tres.

\section{Importación de Malla Gmsh: \texttt{mesh}}
Para problemas continuos 2D, se sustituye el bloque \texttt{nodes}/\texttt{elements} por una referencia al archivo \texttt{.msh}:

\begin{lstlisting}[language=yaml]
mesh: "geometria_placa.msh"
mesh_material: 1
mesh_thickness: 0.1
mesh_quadrature: "2x2"
\end{lstlisting}

\begin{table}[h]
\centering
\begin{tabular}{@{}lp{10cm}@{}}
\toprule
\textbf{Parámetro} & \textbf{Descripción} \\ \midrule
\texttt{mesh} & Ruta al archivo \texttt{.msh} de Gmsh, relativa al YAML. \\
\texttt{mesh\_material} & ID del material por defecto que se asocia a toda la malla. \\
\texttt{mesh\_thickness} & Espesor por defecto (estado plano 2D). \\
\texttt{mesh\_quadrature} & Regla de Gauss: \texttt{"2x2"} (completa, recomendada) o \texttt{"1x1"} (reducida; produce \textit{hourglassing}). \\
\bottomrule
\end{tabular}
\caption{Parámetros del bloque \texttt{mesh}.}
\end{table}

\subsection{Mapeo Avanzado: \texttt{mesh\_physical\_groups}}
Si la geometría tiene zonas con materiales o espesores distintos definidos como \textit{Physical Groups} en Gmsh, se mapean explícitamente:

\begin{lstlisting}[language=yaml]
mesh_physical_groups:
  "Zona_Hormigon":
    material: 1
    thickness: 0.20
    quadrature: "2x2"
  "Zona_Acero":
    material: 2
    thickness: 0.05
\end{lstlisting}

\section{Materiales}
Cada material se declara con un \texttt{id} único y un \texttt{type} registrado en el catálogo:

\begin{lstlisting}[language=yaml]
materials:
  - id: 1
    type: Elastic1D
    E: 200.0e9
  - id: 2
    type: VonMises2D
    E: 210.0e9
    nu: 0.3
    sigma_y: 250.0e6
    H: 2.0e9
    hypothesis: plane_strain
\end{lstlisting}

Los parámetros aceptados dependen del material; consulte el Capítulo~\ref{cap:materiales} para el catálogo completo.

\section{Condiciones de Frontera}
Fenix FEM ofrece tres mecanismos para imponer restricciones cinemáticas (Dirichlet), elegibles según la facilidad operativa:

\subsection{Por nodo: \texttt{boundary\_conditions\_by\_node}}
\begin{lstlisting}[language=yaml]
boundary_conditions_by_node:
  - {node_id: 1, ux: 0.0, uy: 0.0, rz: 0.0}
  - {node_id: 4, ux: 0.0, uy: 0.0}
\end{lstlisting}

\subsection{Por coordenada: \texttt{boundary\_conditions\_by\_coord}}
Escanea todos los nodos y aplica la restricción a aquellos cuya coordenada coincide (dentro de \texttt{tol}):

\begin{lstlisting}[language=yaml]
boundary_conditions_by_coord:
  x_min: {ux: 0.0, uy: 0.0, tol: 1.0e-5}      # empotrar lado izquierdo
  rodillos: {coord: 'y', val: 10.0, uy: 0.0, tol: 1.0e-4}
\end{lstlisting}

Las claves predefinidas \texttt{x\_min}, \texttt{x\_max}, \texttt{y\_min}, \texttt{y\_max} usan los extremos detectados automáticamente del dominio.

\subsection{Por grupo físico: \texttt{boundary\_conditions\_by\_group}}
Aplica restricciones a todos los nodos de un \textit{Physical Group} de Gmsh:

\begin{lstlisting}[language=yaml]
boundary_conditions_by_group:
  "Borde_Empotrado": {ux: 0.0, uy: 0.0}
\end{lstlisting}

\section{Cargas Puntuales}
Misma estructura que las condiciones de frontera, en bloques \texttt{point\_loads\_by\_node}, \texttt{point\_loads\_by\_coord} y \texttt{point\_loads\_by\_group}. Los valores son fuerzas nodales (en las unidades del modelo, típicamente Newtons):

\begin{lstlisting}[language=yaml]
point_loads_by_node:
  - {node_id: 3, uy: -15000.0}

point_loads_by_group:
  "Borde_Carga": {ux: 1500000.0}
\end{lstlisting}

\section{Solver}
\begin{lstlisting}[language=yaml]
solver:
  type: NonlinearSolver
  num_steps: 20
  max_iter: 15
  tol: 1.0e-5
  adaptive: true
\end{lstlisting}

Los solvers disponibles y sus parámetros se documentan en el Capítulo~\ref{cap:solvers}.

\section{Salida: \texttt{output}}
Configura qué se exporta y cuándo:

\begin{lstlisting}[language=yaml]
output:
  file_name: "resultados_placa"
  pre_process: {export: true}     # exporta el modelo sin deformar (paso 0)
  results:
    frequency: "all_steps"        # "last_step" o "all_steps"
    nodal_results: ["Displacements"]
    element_results: ["Von_Mises", "Internal_State", "Sigma_XX"]
  text_export:                     # opcional: salida estilo FEAP en .txt
    export: true
    nodes: all
    elements: all
\end{lstlisting}

\begin{nota}
Cuando \texttt{frequency: "all\_steps"} y el solver es no lineal, se genera un archivo \texttt{.vtu} por paso convergido más un archivo maestro \texttt{.pvd} que ParaView interpreta como animación.
\end{nota}

% ==========================================
\chapter{Catálogo de Elementos Finitos}
\label{cap:elementos}
% ==========================================

El motor expone dos familias de elementos: \textbf{1D} (estructuras reticulares: armaduras, cables, marcos) y \textbf{2D} (continuo plano: cuadrilátero y triángulo). Todos se referencian desde el YAML por su nombre exacto en \texttt{type:}.

\section{Elementos 1D}

\subsection{Armaduras: \texttt{Truss2D} / \texttt{Truss2DCorot} / \texttt{Truss3D} / \texttt{Truss3DCorot}}
Barra biarticulada que transmite exclusivamente esfuerzo axial. Cuatro variantes según dimensión y régimen geométrico:

\begin{table}[h]
\centering
\begin{tabular}{@{}lcccp{4.5cm}@{}}
\toprule
\textbf{Elemento} & \textbf{DOFs/nodo} & \textbf{Dimensión} & \textbf{Régimen geom.} & \textbf{Uso} \\ \midrule
\texttt{Truss2D}      & \texttt{ux,uy}    & 2D & lineal     & cargas pequeñas, rotaciones pequeñas \\
\texttt{Truss2DCorot} & \texttt{ux,uy}    & 2D & corotacional & grandes desplazamientos y rotaciones planas \\
\texttt{Truss3D}      & \texttt{ux,uy,uz} & 3D & lineal     & armaduras espaciales en régimen lineal \\
\texttt{Truss3DCorot} & \texttt{ux,uy,uz} & 3D & corotacional & grandes rotaciones en el espacio \\
\bottomrule
\end{tabular}
\caption{Familia de armaduras.}
\end{table}

\textbf{Parámetros}: \texttt{A} (área de la sección transversal). Convención de signos: $\sigma > 0 \Leftrightarrow \varepsilon > 0$ (elongación).

\textbf{Régimen de validez}: $|\varepsilon| \lesssim 10^{-2}$. Las variantes \texttt{Corot} aceptan rotaciones de cualquier magnitud entre commits; las versiones lineales asumen rotaciones pequeñas.

\begin{lstlisting}[language=yaml]
elements:
  - id: 1
    type: Truss2DCorot
    material: 1
    nodes: [1, 2]
    A: 0.0005
\end{lstlisting}

\subsection{Cables: \texttt{Cable2DCorot} / \texttt{Cable3DCorot}}
Elemento 1D que transmite \emph{únicamente} esfuerzo de tensión (no resiste compresión). La unilateralidad la aporta el material (típicamente \texttt{CableMaterial1D}, ver Capítulo~\ref{cap:materiales}); el elemento implementa la cinemática corotacional y la transferencia de la deformación al material.

\begin{itemize}
    \item \textbf{DOFs/nodo}: \texttt{ux, uy} (2D); \texttt{ux, uy, uz} (3D).
    \item \textbf{Parámetros}: \texttt{A}.
    \item \textbf{Cinemática}: corotacional (longitud y cosenos directores se recalculan en cada evaluación).
    \item \textbf{Régimen de validez}: $|\varepsilon| \lesssim 10^{-2}$ tensado. En estado destensado el elemento aporta rigidez nula al sistema global; otros elementos deben garantizar la estabilidad numérica.
\end{itemize}

\begin{advertencia}
Un cable completamente destensado ($\sigma = 0$) tiene $\mathbf K_T = \mathbf 0$ en sus DOFs. Si todos los caminos de carga del modelo dependen del cable y este se destensa, la matriz tangente global se vuelve singular. Pretensar el cable o garantizar redundancia estructural.
\end{advertencia}

\begin{lstlisting}[language=yaml]
materials:
  - id: 1
    type: CableMaterial1D
    E: 150.0e9

elements:
  - id: 1
    type: Cable2DCorot
    material: 1
    nodes: [1, 2]
    A: 5.0e-5
\end{lstlisting}

\subsection{Marcos 2D: \texttt{Frame2DEuler} / \texttt{Frame2DTimoshenko} / \texttt{Frame2DEulerCorot}}
Elementos viga 2D que transmiten axial, cortante y momento flector. Tres variantes:

\begin{itemize}
    \item \textbf{\texttt{Frame2DEuler}}: vigas esbeltas ($L/h \gtrsim 10$), Euler-Bernoulli, régimen geométricamente lineal. Parámetros: \texttt{A}, \texttt{I}.
    \item \textbf{\texttt{Frame2DTimoshenko}}: vigas peraltadas o cortas ($L/h \lesssim 10$). Incluye deformación por cortante; corrige automáticamente \textit{shear locking} en el límite esbelto. Parámetros: \texttt{A}, \texttt{I}, \texttt{As} (área efectiva de cortante), \texttt{nu} (opcional si el material expone \texttt{nu}).
    \item \textbf{\texttt{Frame2DEulerCorot}}: variante corotacional de Euler-Bernoulli. Captura grandes desplazamientos y grandes rotaciones rígidas del elemento (rotaciones deformacionales nodales moderadas, $|\bar\theta| \lesssim 30^\circ$). Parámetros: \texttt{A}, \texttt{I}.
\end{itemize}

DOFs/nodo: \texttt{ux, uy, rz}. Convención: $\sigma > 0 \Leftrightarrow \varepsilon > 0$ (elongación) en el eje axial; $r_z$ positivo antihorario.

\begin{advertencia}
\textbf{Plasticidad por flexión.} Los elementos marco pasan al material únicamente la deformación axial centroidal $(u_j - u_i)/L$; el módulo tangente $E_t$ devuelto escala toda la matriz local (axial y flexión por igual). Esto significa que los marcos reproducen correctamente \emph{plasticidad axial pura} pero \emph{no} forman rótulas plásticas por flexión: la fluencia por momento requiere integración de $\sigma(y)$ sobre la sección, característica que se incorporará en una clase \texttt{FiberSection} futura.
\end{advertencia}

\begin{lstlisting}[language=yaml]
elements:
  - id: 1
    type: Frame2DEulerCorot
    material: 1
    nodes: [1, 2]
    A: 0.01
    I: 8.33e-6

  - id: 2
    type: Frame2DTimoshenko
    material: 1
    nodes: [2, 3]
    A: 0.01
    I: 8.33e-6
    As: 0.00833
\end{lstlisting}

\subsection{Marco 3D: \texttt{Frame3D}}
Viga 3D Euler-Bernoulli con 6 DOFs/nodo (\texttt{ux, uy, uz, rx, ry, rz}). Transmite axial, cortantes en dos planos, flexiones en dos planos y torsión de Saint-Venant pura. Régimen geométricamente lineal (la variante corotacional 3D queda como pendiente).

\textbf{Parámetros}:
\begin{itemize}
    \item \texttt{A} --- área de la sección.
    \item \texttt{Iy}, \texttt{Iz} --- momentos de inercia respecto a los ejes locales $y$, $z$.
    \item \texttt{J} --- constante torsional de Saint-Venant.
    \item \texttt{nu} --- coeficiente de Poisson (opcional; se toma del material si lo expone).
    \item \texttt{ref\_vector} --- vector de referencia (opcional, default \texttt{[0,0,1]}) que fija la orientación de los ejes locales $y$, $z$ de la sección. Para barras casi-verticales se sugiere fijarlo explícitamente.
\end{itemize}

\begin{lstlisting}[language=yaml]
elements:
  - id: 1
    type: Frame3D
    material: 1
    nodes: [1, 2]
    A: 0.01
    Iy: 8.33e-6
    Iz: 8.33e-6
    J: 1.66e-5
    ref_vector: [0.0, 0.0, 1.0]
\end{lstlisting}

\section{Elementos 2D}

La importación de mallas desde Gmsh instancia automáticamente los elementos 2D según la topología detectada (cuadriláteros $\to$ \texttt{Quad4}; triángulos $\to$ \texttt{Tri3}). También pueden declararse manualmente desde el bloque \texttt{elements}.

\subsection{Cuadrilátero Bilineal: \texttt{Quad4}}
Elemento isoparamétrico de 4 nodos, integración Gauss configurable.

\begin{itemize}
    \item \textbf{DOFs/nodo}: \texttt{ux, uy} \quad ($\texttt{STRAIN\_DIM} = 3$, Voigt $[\varepsilon_{xx}, \varepsilon_{yy}, \gamma_{xy}]$).
    \item \textbf{Nodos}: 4, ordenados en sentido antihorario.
    \item \textbf{Cuadratura}: \texttt{"2x2"} por defecto (4 puntos). \texttt{"1x1"} disponible pero produce modos espurios (\textit{hourglassing}).
    \item \textbf{Parámetros}: \texttt{thickness}, \texttt{quadrature} (opcional).
    \item \textbf{Implementación}: kernels críticos compilados con \texttt{@njit} (Numba).
    \item \textbf{Limitación}: bloqueo volumétrico con materiales casi-incompresibles ($\nu \to 0.5$); en ese régimen requeriría formulación mixta (no implementada).
\end{itemize}

\subsection{Triángulo de Deformación Constante: \texttt{Tri3}}
Triángulo CST de 3 nodos, un único punto de integración.

\begin{itemize}
    \item \textbf{DOFs/nodo}: \texttt{ux, uy} \quad ($\texttt{STRAIN\_DIM} = 3$).
    \item \textbf{Cuadratura}: 1 punto (deformación uniforme).
    \item \textbf{Parámetros}: \texttt{thickness}.
    \item \textbf{Limitación}: \textit{shear locking} severo, convergencia lenta. \textbf{Preferir \texttt{Quad4}} salvo en transiciones donde Quad4 no encaja geométricamente.
\end{itemize}

% ==========================================
\chapter{Catálogo de Modelos Constitutivos}
\label{cap:materiales}
% ==========================================

\section{Elasticidad Lineal}

\subsection{\texttt{Elastic1D} --- elasticidad lineal axial}
Ley de Hooke escalar $\sigma = E \cdot \varepsilon$. Sin variables internas. Compatible con todos los elementos 1D (armaduras y marcos).

\begin{lstlisting}[language=yaml]
materials:
  - id: 1
    type: Elastic1D
    E: 200.0e9
\end{lstlisting}

\subsection{\texttt{Elastic2D} --- elasticidad lineal isótropa 2D}
Tensor constitutivo isotrópico $\boldsymbol\sigma = \mathbf C \cdot \boldsymbol\varepsilon$ en notación Voigt $[xx, yy, xy]$. Sin variables internas. Compatible con \texttt{Quad4}, \texttt{Tri3}.

\begin{lstlisting}[language=yaml]
materials:
  - id: 1
    type: Elastic2D
    E: 200.0e9
    nu: 0.3
    hypothesis: plane_stress
\end{lstlisting}

La hipótesis \texttt{plane\_stress} se usa para placas delgadas sin confinamiento normal; \texttt{plane\_strain} para secciones de presa, túneles o problemas con extensión infinita en el eje longitudinal.

\section{Plasticidad}

\subsection{\texttt{Elastoplastic1D} --- plasticidad J2 axial con endurecimiento isotrópico}
Plasticidad asociativa con criterio $f = |\sigma_{\text{trial}}| - (\sigma_y + H \alpha) \le 0$. Algoritmo de \textit{return mapping} clásico 1D; tangente algorítmica consistente $E_t = E\,H / (E + H)$.

\begin{itemize}
    \item \textbf{Parámetros}: \texttt{E}, \texttt{sigma\_y} (fluencia inicial), \texttt{H} (módulo de endurecimiento; \texttt{H=0} = perfectamente plástico).
    \item \textbf{Variables internas}: $\varepsilon_p$ (deformación plástica), $\alpha$ (acumulada equivalente).
    \item \textbf{Compatible con}: armaduras y marcos 1D (en marcos solo se aplica al esfuerzo axial; ver advertencia en el Capítulo~\ref{cap:elementos}).
\end{itemize}

\begin{lstlisting}[language=yaml]
materials:
  - id: 1
    type: Elastoplastic1D
    E: 200.0e9
    sigma_y: 250.0e6
    H: 2.0e9
\end{lstlisting}

\subsection{\texttt{VonMises2D} --- plasticidad J2 plana con endurecimiento isotrópico}
Return mapping radial sobre la parte desviadora. Tangente algorítmica consistente.

\begin{itemize}
    \item \textbf{Parámetros}: \texttt{E}, \texttt{nu}, \texttt{sigma\_y}, \texttt{H}, \texttt{hypothesis} (\textbf{obligatorio} \texttt{plane\_strain}).
    \item \textbf{Variables internas}: $\varepsilon_p$ tensorial (4 componentes $[xx, yy, zz, xy]$), $\alpha$.
    \item \textbf{Implementación}: kernel \texttt{\_compute\_j2\_plasticity} con \texttt{@njit}.
\end{itemize}

\begin{advertencia}
\textbf{Restricción cinemática}: solo \texttt{plane\_strain}. Asignar \texttt{plane\_stress} a \texttt{VonMises2D} lanza \texttt{NotImplementedError} (el return mapping en estado plano de esfuerzos requiere algoritmo distinto, no implementado).
\end{advertencia}

\begin{lstlisting}[language=yaml]
materials:
  - id: 2
    type: VonMises2D
    E: 210.0e9
    nu: 0.3
    sigma_y: 250.0e6
    H: 2.0e9
    hypothesis: plane_strain
\end{lstlisting}

\section{Daño Mecánico}

\subsection{\texttt{IsotropicDamage1D} --- daño isótropo 1D con softening exponencial}
Daño escalar $d \in [0, 1)$ con módulo secante $E_{\text{sec}} = (1 - d)\,E$. Evolución por norma de la deformación equivalente $\varepsilon_{eq} = |\varepsilon|$ y umbral $\kappa_0$:

$$d = 1 - \frac{\kappa_0}{\kappa}\,\exp\bigl(-\alpha(\kappa - \kappa_0)\bigr) \quad \text{para } \kappa > \kappa_0$$

\begin{itemize}
    \item \textbf{Parámetros}: \texttt{E}, \texttt{kappa\_0} (umbral elástico), \texttt{alpha} (velocidad de softening).
    \item \textbf{Variables internas}: $\kappa$ (máxima histórica), $d$ (daño).
    \item \textbf{Tangente}: secante (no algorítmica consistente; adecuado para Newton amortiguado, no óptimo para arc-length agresivo).
\end{itemize}

\subsection{\texttt{IsotropicDamage2D} --- daño isótropo 2D}
Extensión 2D del modelo anterior. Tensor constitutivo secante $\mathbf C_{\text{sec}} = (1-d)\,\mathbf C_e$.

\begin{itemize}
    \item \textbf{Deformación equivalente}: $\varepsilon_{eq} = \sqrt{\varepsilon_{xx}^2 + \varepsilon_{yy}^2 + \tfrac{1}{2}\gamma_{xy}^2}$.
    \item \textbf{Parámetros}: \texttt{E}, \texttt{nu}, \texttt{kappa\_0}, \texttt{alpha}, \texttt{hypothesis}.
    \item \textbf{Limitación}: $\varepsilon_{eq}$ simétrica no distingue daño en tensión vs compresión; para hormigón se requiere modelo de Mazars o split tensión/compresión (no implementado).
\end{itemize}

\begin{lstlisting}[language=yaml]
materials:
  - id: 3
    type: IsotropicDamage2D
    E: 30.0e9
    nu: 0.2
    kappa_0: 1.5e-4
    alpha: 800.0
    hypothesis: plane_stress
\end{lstlisting}

\section{Elasticidad Unilateral (Cable)}

\subsection{\texttt{CableMaterial1D} --- cable lineal en tensión, nulo en compresión}
Material 1D \emph{memoryless} que captura la propiedad esencial del cable: solo transmite esfuerzo cuando está tensado.

$$\sigma(\varepsilon) = E\,\langle\varepsilon\rangle^+ = \begin{cases} E\,\varepsilon & \varepsilon > 0 \\ 0 & \varepsilon \le 0 \end{cases}$$

con $\langle \cdot \rangle^+$ el operador de MacAuley (parte positiva). Tangente discontinua en $\varepsilon = 0$ (asignación al régimen compresivo por convención).

\begin{itemize}
    \item \textbf{Parámetros}: \texttt{E} (módulo de Young en tensión).
    \item \textbf{Sin variables internas} (la respuesta depende sólo de $\varepsilon$ instantánea).
    \item \textbf{Caveats numéricos}: en transiciones tensado $\leftrightarrow$ destensado Newton-Raphson puede oscilar. Mitigación: usar \texttt{ArcLengthSolver} o pasos finos con \texttt{adaptive}.
\end{itemize}

\begin{lstlisting}[language=yaml]
materials:
  - id: 1
    type: CableMaterial1D
    E: 150.0e9
\end{lstlisting}

% ==========================================
\chapter{Esquemas de Solución}
\label{cap:solvers}
% ==========================================

\section{\texttt{LinearSolver} --- solución directa lineal}
Resuelve $\mathbf K \cdot \mathbf U = \mathbf F$ en un único \texttt{spsolve}. Aplicación de condiciones de frontera por penalidad.

\begin{itemize}
    \item \textbf{Cuándo usarlo}: problemas estrictamente lineales (todos los materiales con tangente constante, sin grandes desplazamientos, sin contacto).
    \item \textbf{Parámetros}: \texttt{penalty\_value} (default \texttt{1.0e15}).
    \item \textbf{No aplica}: cualquier no-linealidad material (daño, plasticidad, cable) o geométrica (corotacional). Producirá resultados inconsistentes.
\end{itemize}

\begin{lstlisting}[language=yaml]
solver:
  type: LinearSolver
\end{lstlisting}

\section{\texttt{NonlinearSolver} --- Newton-Raphson incremental con paso adaptativo}
Bucle externo de incrementos del factor de carga $\lambda \in [0, 1]$; bucle interno Newton-Raphson sobre $\mathbf K_t \cdot \Delta\mathbf U = \mathbf R = \lambda\,\mathbf F_{\text{ext}} - \mathbf F_{\text{int}}$.

\textbf{Criterio de convergencia dual}:
$$\max(\text{err}_{\text{disp}}, \text{err}_{\text{force}}) < \text{tol}$$
con $\text{err}_{\text{disp}} = \lVert\Delta\mathbf U\rVert / (\lVert\mathbf U\rVert + \epsilon)$ y $\text{err}_{\text{force}} = \lVert\mathbf R\rVert / \max(\lVert\lambda\,\mathbf F_{\text{ext}}\rVert, \lVert\mathbf F_{\text{int}}\rVert, \epsilon)$.

\textbf{Adaptatividad} (\texttt{adaptive: true}):
\begin{itemize}
    \item Si converge en menos de 5 iteraciones, agranda el siguiente paso ($\times 1.5$).
    \item Si no converge, biseca el paso ($\div 2$). Falla definitivamente si $\Delta\lambda < \text{min\_delta\_lambda}$.
\end{itemize}

\begin{lstlisting}[language=yaml]
solver:
  type: NonlinearSolver
  num_steps: 20
  max_iter: 15
  tol: 1.0e-5
  adaptive: true
\end{lstlisting}

\textbf{Parámetros}: \texttt{tol}, \texttt{max\_iter}, \texttt{num\_steps}, \texttt{adaptive}, \texttt{penalty\_value}, \texttt{min\_delta\_lambda}.

\textbf{Limitación}: no puede atravesar puntos límite (snap-through) ni recorrer ramas con derivada negativa de la curva carga-desplazamiento. Para esos casos $\to$ \texttt{ArcLengthSolver}.

\section{\texttt{ArcLengthSolver} --- longitud de arco cilíndrico (Crisfield)}
Traza curvas de equilibrio con snap-through, snap-back o pérdida de unicidad de carga, controlando simultáneamente desplazamientos y factor de carga $\lambda$.

\textbf{Esquema}:
\begin{itemize}
    \item \textbf{Predictor tangente}: $\Delta\mathbf U_t = \mathbf K_t^{-1} \cdot \mathbf F_{\text{ext}}^{\text{ref}}$; $\Delta\lambda = \pm\,dl / \lVert\Delta\mathbf U_t\rVert$. Signo elegido por proyección con el incremento previo.
    \item \textbf{Corrector iterativo}: en cada iteración resuelve dos sistemas ($\mathbf K \cdot \Delta\mathbf U_R = \mathbf R$ y $\mathbf K \cdot \Delta\mathbf U_t = \mathbf F_{\text{ext}}$) y aplica la restricción cuadrática cilíndrica $\lVert\delta\mathbf U\rVert^2 = dl^2$.
    \item \textbf{Auto-ajuste de \texttt{dl}}: se agranda en convergencias rápidas, se reduce en lentas, biseca al fallar.
    \item \textbf{Caso especial \texttt{final\_step}}: si el predictor sobrepasaría \texttt{max\_lambda}, fija $\lambda$ a ese valor y resuelve solo desplazamientos (Newton-Raphson puro).
\end{itemize}

\begin{lstlisting}[language=yaml]
solver:
  type: ArcLengthSolver
  tol: 1.0e-5
  max_iter: 15
  max_lambda: 1.0
  initial_dl: 0.05
  max_steps: 200
\end{lstlisting}

\textbf{Cuándo usarlo}: problemas con softening pronunciado (daño, post-pandeo, snap-through de cúpulas), o cuando \texttt{NonlinearSolver} diverge cerca de un punto límite. Más caro por paso (dos \texttt{spsolve} por iteración) pero indispensable en estos casos.

\textbf{Referencia}: Crisfield, ``A fast incremental/iterative solution procedure that handles snap-through'' (\textit{Computers \& Structures}, 1981).

% ==========================================
\chapter{Post-Procesamiento y Visualización}
% ==========================================

\section{Formato VTU (VTK XML)}
Mediante \texttt{meshio}, Fenix FEM exporta los resultados en archivos \texttt{.vtu} compatibles con ParaView. La información se clasifica en dos diccionarios estándar:

\subsection{Datos Nodales (\textit{Point Data})}
\begin{itemize}
    \item \texttt{Displacements}: campo vectorial $(u_x, u_y, u_z)$ por nodo.
    \item \texttt{External\_Forces}: distribución nodal del vector de cargas aplicado en el paso actual.
    \item \texttt{Supports}: indicador booleano de DOFs restringidos.
\end{itemize}

\subsection{Datos de Elemento (\textit{Cell Data})}
\begin{itemize}
    \item \texttt{Von\_Mises}: tensión equivalente J2 (cuando aplica).
    \item \texttt{Internal\_State}: variable principal del material según su \texttt{PRIMARY\_STATE\_VAR} (acumulada equivalente $\alpha$ en plasticidad, $d$ en daño, etc.).
    \item \texttt{Sigma\_XX}, \texttt{Sigma\_YY}, \texttt{Tau\_XY}: componentes del tensor de esfuerzos en notación Voigt.
\end{itemize}

\section{Animación con Archivo PVD}
Cuando se exporta con \texttt{frequency: "all\_steps"}, además de un \texttt{.vtu} por paso se genera un archivo maestro \texttt{.pvd} (XML colección) que ParaView abre como animación temporal. El ``tiempo'' en el archivo PVD coincide con el factor de carga $\lambda$.

\section{Salida en Texto Plano (Estilo FEAP)}
Opcional para verificación manual o post-procesamiento con scripts. Configurable por nodos y elementos seleccionados:

\begin{lstlisting}[language=yaml]
output:
  text_export:
    export: true
    nodes: [1, 5, 10]                # o "all"
    elements: all
    nodal_results: ["Displacements", "External_Forces"]
    element_results: ["Von_Mises", "Internal_State"]
\end{lstlisting}

\section{Workflow ParaView}
\begin{enumerate}
    \item \texttt{File > Open} y seleccionar el archivo \texttt{.pvd} (o el grupo de \texttt{.vtu}).
    \item \textit{Apply} en el panel \textit{Properties}.
    \item Reproducir la animación con los controles temporales de la barra superior.
    \item Aplicar filtro \textbf{Warp By Vector} con \texttt{Displacements} para visualizar la deformación física (ajustar \texttt{Scale Factor} para amplificar).
    \item Cambiar el coloreado de \texttt{Solid Color} a la variable de interés (\texttt{Von\_Mises}, \texttt{Internal\_State}, \texttt{Sigma\_XX}, \ldots).
\end{enumerate}

% ==========================================
\chapter{Ejemplos Completos}
% ==========================================

\section{Marco 2D --- Viga en Voladizo}
Combinación de un tramo Euler-Bernoulli y un tramo Timoshenko, empotrada en el extremo y cargada en la punta.

\begin{lstlisting}[language=yaml, caption=Marco 2D mixto Euler/Timoshenko]
nodes:
  - {id: 1, coords: [0.0, 0.0]}
  - {id: 2, coords: [2.0, 0.0]}
  - {id: 3, coords: [4.0, 0.0]}

materials:
  - id: 1
    type: Elastic1D
    E: 200.0e9

elements:
  - id: 1
    type: Frame2DEuler
    material: 1
    nodes: [1, 2]
    A: 0.01
    I: 8.33e-6

  - id: 2
    type: Frame2DTimoshenko
    material: 1
    nodes: [2, 3]
    A: 0.01
    I: 8.33e-6
    As: 0.00833

boundary_conditions_by_node:
  - {node_id: 1, ux: 0.0, uy: 0.0, rz: 0.0}

point_loads_by_node:
  - {node_id: 3, uy: -15000.0}

solver:
  type: LinearSolver

output:
  file_name: "viga_voladizo_marco"
  pre_process: {export: true}
  results:
    frequency: "last_step"
    nodal_results: ["Displacements"]
\end{lstlisting}

\section{Placa Continua con Plasticidad J2 (Gmsh)}
Placa con perforación central importada desde \texttt{.msh}, sometida a carga axial creciente hasta superar el límite elástico del acero. Solver no lineal con paso adaptativo.

\begin{lstlisting}[language=yaml, caption=Placa perforada elastoplástica]
mesh: "placa_agujero.msh"

mesh_physical_groups:
  "Superficie_Placa":
    material: 1
    thickness: 0.05
    quadrature: "2x2"

materials:
  - id: 1
    type: VonMises2D
    E: 200.0e9
    nu: 0.3
    sigma_y: 250.0e6
    H: 1.5e9
    hypothesis: plane_strain

boundary_conditions_by_group:
  "Borde_Fijo": {ux: 0.0, uy: 0.0}

point_loads_by_group:
  "Borde_Carga": {ux: 1500000.0}

solver:
  type: NonlinearSolver
  num_steps: 30
  max_iter: 10
  tol: 1.0e-5
  adaptive: true

output:
  file_name: "estudio_placa"
  pre_process: {export: true}
  results:
    frequency: "all_steps"
    nodal_results: ["Displacements"]
    element_results: ["Von_Mises", "Internal_State"]
\end{lstlisting}

\textbf{Interpretación}: durante los primeros pasos (régimen elástico) Newton-Raphson converge en 1--2 iteraciones. Cuando la concentración de esfuerzos en la periferia del orificio supera $\sigma_y = 250$ MPa, comienzan iteraciones del \textit{return mapping}; \texttt{Internal\_State} (la deformación plástica acumulada $\alpha$) crece visiblemente alrededor del agujero.

\section{Cable Pretensado en Régimen Corotacional}
Cable 2D con material unilateral, cargado transversalmente en su nodo central. Demuestra el acoplamiento elemento corotacional + material unilateral.

\begin{lstlisting}[language=yaml, caption=Cable bajo carga transversal]
nodes:
  - {id: 1, coords: [0.0, 0.0]}
  - {id: 2, coords: [5.0, 0.0]}
  - {id: 3, coords: [10.0, 0.0]}

materials:
  - id: 1
    type: CableMaterial1D
    E: 150.0e9

elements:
  - id: 1
    type: Cable2DCorot
    material: 1
    nodes: [1, 2]
    A: 5.0e-5

  - id: 2
    type: Cable2DCorot
    material: 1
    nodes: [2, 3]
    A: 5.0e-5

boundary_conditions_by_node:
  - {node_id: 1, ux: 0.0, uy: 0.0}
  - {node_id: 3, ux: 0.0, uy: 0.0}

point_loads_by_node:
  - {node_id: 2, uy: -500.0}

solver:
  type: NonlinearSolver
  num_steps: 50
  max_iter: 25
  tol: 1.0e-5
  adaptive: true

output:
  file_name: "cable_carga_transversal"
  pre_process: {export: true}
  results:
    frequency: "all_steps"
    nodal_results: ["Displacements"]
\end{lstlisting}

% ==========================================
\chapter{Uso Avanzado: API Programática}
% ==========================================

Aunque el flujo principal es \textit{data-driven} (YAML), el motor se puede usar como librería Python para escenarios donde el archivo no es suficiente: optimización paramétrica, generación procedural de mallas, acoplamiento con bibliotecas externas, scripting de batería de tests.

\section{Patrón Fachada}
Los componentes principales se exponen desde la raíz del paquete:

\begin{lstlisting}[language=Python, caption=Script estructural directo]
import numpy as np
from fenix import Domain, VonMises2D, Quad4, Assembler, NonlinearSolver, VtkExporter

# 1. Dominio y nodos
domain = Domain()
n1 = domain.add_node(1, [0.0, 0.0])
n2 = domain.add_node(2, [1.0, 0.0])
n3 = domain.add_node(3, [1.0, 1.0])
n4 = domain.add_node(4, [0.0, 1.0])

# 2. Material y elemento
acero = VonMises2D(E=200e9, nu=0.3, sigma_y=250e6, H=2e9, hypothesis='plane_strain')
placa = Quad4(element_id=1, nodes=[n1, n2, n3, n4], material=acero, thickness=0.1)
domain.add_element(placa)

# 3. Restricciones (Dirichlet)
n1.fix_dof('ux', 0.0); n1.fix_dof('uy', 0.0)
n4.fix_dof('ux', 0.0); n4.fix_dof('uy', 0.0)

domain.generate_equation_numbers()

# 4. Cargas (Neumann)
F_ext = np.zeros(domain.total_dofs)
F_ext[n2.dofs['ux']] = 150000.0
F_ext[n3.dofs['ux']] = 150000.0

# 5. Resolución
assembler = Assembler(domain)
solver = NonlinearSolver(assembler, num_steps=10, tol=1e-5, adaptive=True)
U = solver.solve(F_ext)

# 6. Exportación
VtkExporter(domain).export("resultados.vtu", U=U, F_ext=F_ext)
\end{lstlisting}

% ==========================================
\chapter{Diagnóstico de Problemas}
% ==========================================

\begin{itemize}
    \item \textbf{\textit{YamlValidationError: parámetro `X' no aceptado por `Y'.}} \\
    El parser detecta un parámetro que el constructor del elemento no admite. La salida lista los parámetros aceptados; revisar el catálogo o la spec del componente.

    \item \textbf{\textit{Matriz singular detectada.}} \\
    Restricciones cinemáticas insuficientes: el dominio puede trasladarse o rotar como cuerpo rígido. Verificar que al menos un nodo tiene fijos los suficientes DOFs para anclar el dominio. En modelos con cables o materiales con softening pronunciado, comprobar que existen caminos de carga alternativos al elemento que se destensa o degrada.

    \item \textbf{\textit{Newton-Raphson no convergió / bisecando incremento.}} \\
    El paso de carga es demasiado grande. Aumentar \texttt{num\_steps}, asegurar \texttt{adaptive: true}. Verificar tipeo de los parámetros constitutivos ($\sigma_y$, $\kappa_0$, etc.). Si el problema tiene snap-through o softening, cambiar a \texttt{ArcLengthSolver}.

    \item \textbf{\textit{No se importó ningún elemento (Solid2D).}} \\
    El parser de Gmsh no encontró superficies bidimensionales. En Gmsh, crear una \textit{Plane Surface} y generar la malla 2D antes de exportar.

    \item \textbf{\textit{NotImplementedError: VonMises2D requires plane\_strain.}} \\
    El modelo \texttt{VonMises2D} solo soporta deformación plana. Cambiar \texttt{hypothesis} a \texttt{plane\_strain} en el material o reemplazar por un modelo compatible con plane stress (no implementado actualmente).

    \item \textbf{\textit{Cable totalmente destensado --- $\mathbf K_T = \mathbf 0$.}} \\
    Pretensar el cable o garantizar que la estructura tiene rigidez residual por otros elementos. Considerar también pasos de carga finos para capturar la transición tensado $\to$ destensado.
\end{itemize}

\newpage
\phantomsection
\addcontentsline{toc}{chapter}{Índice Alfabético}
\renewcommand{\indexname}{Índice Alfabético}
\printindex

\end{document}
